\section{Introduction}
\label{sec:introduction}

\subsection{Purpose of this report}

This document consolidates the working history of RoboReason ROS2 into a
single reference, written at a natural pause point (a planned one-month
break before resuming in September 2026). Its two goals are:

\begin{enumerate}[nosep]
  \item Preserve \emph{why} decisions were made, not just what the code
    currently does -- the git log and the code itself already answer
    ``what''; this report exists mainly to answer ``why'', which decays
    fastest from memory.
  \item Give a concrete, prioritized starting point for the next working
    session, so the first hours back are spent on the highest-value item
    rather than re-deriving context.
\end{enumerate}

Two companion documents remain the canonical low-level references and are
not duplicated here in full: \texttt{docs/SESSION\_CONTEXT.md} (a detailed,
append-only session log with exact file-by-file change tables) and
\texttt{docs/TODO.md} (the raw backlog). This report synthesizes both into a
narrative and adds the benchmark study, timeline, and roadmap framing that
neither file provides on its own.

\subsection{Project goal}

Given a natural-language instruction and a workspace containing a small set
of known object types (currently four colored cubes and a tray, on a table
within reach of a UR5cb arm), produce and execute a safe, correct sequence of
pick-and-place actions -- without hand-coding task-specific logic per
instruction. The system should degrade gracefully across three axes of task
difficulty that recur throughout this report and the benchmark study:

\begin{itemize}[nosep]
  \item \textbf{Task length} -- a single pick-and-place vs.\ a multi-object
    sequence.
  \item \textbf{Specificity} -- referring to an object directly (``the red
    cube'') vs.\ by a shared, ambiguous class (``the cubes'') or via a
    computed reference (an arithmetic condition on a stated value mapping).
  \item \textbf{Affordance / perception load} -- whether object positions are
    already known (LLM mode, static scene) or must be perceived from a live
    camera frame (VLM mode).
\end{itemize}

\subsection{From VLA to LLM/VLM: the architectural pivot}

Documents predating this repository's git history (found in the sibling
\texttt{robotic-ai/docs} directory: \emph{Robot AI -- Technical Steps},
\emph{Vision-Language-Action (VLA) Models in Robotics}, \emph{VLA Future
Directions}) describe an earlier plan built around end-to-end
Vision-Language-Action models. That direction was abandoned before the
current repository's initial commit. It is documented here for completeness
and explicitly \textbf{not} as a currently accurate description of the
system -- those files should not be consulted for architecture questions
about the present codebase.

The system that actually exists is a classical planning pipeline with
foundation models substituting for a hand-written planner and, in VLM mode, a
hand-written perception module:

\begin{center}
\begin{tabular}{@{}p{3.2cm}p{5.6cm}p{5.6cm}@{}}
\toprule
\textbf{Layer} & \textbf{VLA approach (abandoned)} & \textbf{Current LLM/VLM approach} \\
\midrule
Perception $\to$ action & Single end-to-end network, pixels + language $\to$
  low-level action & Foundation VLM $\to$ structured JSON plan of named
  skills; skills executed by hand-written trajectory code \\
Training & Requires task-specific demonstration data and fine-tuning &
  Zero-shot / prompted use of off-the-shelf hosted models \\
Debuggability & Opaque -- a wrong action has no intermediate artifact to
  inspect & Every plan step, pixel grounding, and deprojection is logged to
  \texttt{debug/<run\_id>/} and independently inspectable \\
Geometry correctness & Learned implicitly from data & Computed
  deterministically in Python wherever the model's own arithmetic proved
  unreliable (Section~\ref{sec:tech-decisions}) \\
\bottomrule
\end{tabular}
\end{center}

\subsection{Physical setup}

\begin{itemize}[nosep]
  \item \textbf{Arm}: Universal Robots UR5cb, controlled via the standard
    \texttt{ur\_robot\_driver} action/controller interface.
  \item \textbf{Gripper}: OnRobot RG2, driven via digital I/O (URScript
    socket control was tried and replaced early in the project history, see
    commit \texttt{92d0ead}).
  \item \textbf{Camera}: Orbbec RGB-D camera, used for VLM-mode scene capture
    and ChArUco-based extrinsic calibration (board-to-base transform).
  \item \textbf{Workspace}: a table within the arm's reach, a tray as a
    placement target, and (as of the current physical setup, changed during
    Session 5) four cubes of identical shape distinguished only by color --
    blue, red, white, orange.
\end{itemize}
